\paragraph{Prior tools from Section 2.1.}
Our proof of Theorem 1.2 is based on the meta-complexity of time-bounded Kolmogorov complexity: We crucially use the fact that time-bounded Kolmogorov complexity can be approximated in Heuristica.
\paragraph{Lemma 2.1 ([31,33]).}
If $\mathsf{DistNP}\subseteq\mathsf{AvgP}$, then $\mathsf{GapMINKT}\in\mathsf P$.
The second lemma is the reconstruction property of a $k$-wise direct product generator [34], which turned out to be a fundamental tool for analyzing Kolmogorov complexity [33,65,32,35]. A $k$-wise direct product generator $\mathsf{DP}_k:\{0,1\}^n\times\{0,1\}^{nk}\to\{0,1\}^{nk+k}$ is defined as follows:
\[\mathsf{DP}_k(x;z):=(z^1,\ldots,z^k,\langle z^1,x\rangle,\ldots,\langle z^k,x\rangle)\]
for every $x\in\{0,1\}^n$ and every $z=(z^1,\ldots,z^k)\in(\{0,1\}^n)^k$, where $\langle x,y\rangle$ denotes the inner product of $x$ and $y\in\{0,1\}^n$ over $\mathrm{GF}(2)$, i.e., $\langle x,y\rangle:=(\sum_{i=1}^nx_iy_i)\bmod2$. The $k$-wise direct product generator $\mathsf{DP}_k(x;-)$ is a pseudorandom generator secure against an algorithm $D$ if $K(x\mid D)>k+O(\log|x|)$. More formally, we have the following property:
\paragraph{Lemma 2.2 (Deterministic Reconstruction for $\mathsf{DP}_k$; see [35]).}
Assume that $\mathsf E\not\subseteq\mathsf{i.o.SIZE}(2^{\epsilon n})$ for some constant $\epsilon>0$. Then, there exists a polynomial $p$ such that, for every $n\in\mathbb N$, $x\in\{0,1\}^n$, parameters $k,\epsilon^{-1},s\in\mathbb N$, and for every randomized circuit $D$ of size $s$ such that
\[\left|\Pr_{z,r}[D(\mathsf{DP}_k(x;z);r)=1]-\Pr_{w,r}[D(w;r)=1]\right|\geq\epsilon,\]
where $z\sim\{0,1\}^{nk}$, $w\sim\{0,1\}^{nk+k}$, and $r\sim\{0,1\}^s$, it holds that
\[K^{p(ns/\epsilon)}(x\mid D)\leq k+\log p(ns/\epsilon).\]
The assumption of Lemma 2.2 is satisfied in Heuristica, as Buhrman et al. [18] showed that $\mathsf E\not\subseteq\mathsf{i.o.SIZE}(2^{\epsilon n})$ for some constant $\epsilon>0$ if $\mathsf{DistNP}\subseteq\mathsf{AvgP}$. We mention in passing that the proof of Lemma 2.1 also uses Lemma 2.2.
